\documentclass[10pt,a4paper]{amsart}
\usepackage{amsmath,amssymb,amsthm,graphicx}
\usepackage{verbatim}
\newenvironment{reference}{\comment}{\endcomment}
%\newenvironment{reference}{}{}
\newenvironment{slogan}{\comment}{\endcomment}
\newenvironment{history}{\comment}{\endcomment}

% For commutative diagrams we use Xy-pic
\usepackage[all]{xy}

% We use 2cell for 2-commutative diagrams.
\xyoption{2cell}
% Theorem environments.
%
\theoremstyle{plain}
\newtheorem{theorem}[subsection]{Theorem}
\newtheorem{proposition}[subsection]{Proposition}
\newtheorem{lemma}[subsection]{Lemma}

\theoremstyle{definition}
\newtheorem{definition}[subsection]{Definition}
\newtheorem{example}[subsection]{Example}
\newtheorem{exercise}[subsection]{Exercise}
\newtheorem{situation}[subsection]{Situation}

\theoremstyle{remark}
\newtheorem{remark}[subsection]{Remark}
\newtheorem{remarks}[subsection]{Remarks}

\numberwithin{equation}{subsection}

% Macros
%
\def\lim{\mathop{\mathrm{lim}}\nolimits}
\def\colim{\mathop{\mathrm{colim}}\nolimits}
\def\Spec{\mathop{\mathrm{Spec}}}
\def\Hom{\mathop{\mathrm{Hom}}\nolimits}
\def\Ext{\mathop{\mathrm{Ext}}\nolimits}
\def\SheafHom{\mathop{\mathcal{H}\!\mathit{om}}\nolimits}
\def\SheafExt{\mathop{\mathcal{E}\!\mathit{xt}}\nolimits}
\def\Sch{\mathit{Sch}}
\def\Mor{\mathop{\mathrm{Mor}}\nolimits}
\def\Ob{\mathop{\mathrm{Ob}}\nolimits}
\def\Sh{\mathop{\mathit{Sh}}\nolimits}
\def\NL{\mathop{N\!L}\nolimits}
\def\CH{\mathop{\mathrm{CH}}\nolimits}
\def\proetale{{pro\text{-}\acute{e}tale}}
\def\etale{{\acute{e}tale}}
\def\QCoh{\mathit{QCoh}}
\def\Ker{\mathop{\mathrm{Ker}}}
\def\Im{\mathop{\mathrm{Im}}}
\def\Coker{\mathop{\mathrm{Coker}}}
\def\Coim{\mathop{\mathrm{Coim}}}

% Boxtimes
%
\DeclareMathSymbol{\boxtimes}{\mathbin}{AMSa}{"02}

%
% Macros for moduli stacks/spaces
%
\def\QCohstack{\mathcal{QC}\!\mathit{oh}}
\def\Cohstack{\mathcal{C}\!\mathit{oh}}
\def\Spacesstack{\mathcal{S}\!\mathit{paces}}
\def\Quotfunctor{\mathrm{Quot}}
\def\Hilbfunctor{\mathrm{Hilb}}
\def\Curvesstack{\mathcal{C}\!\mathit{urves}}
\def\Polarizedstack{\mathcal{P}\!\mathit{olarized}}
\def\Complexesstack{\mathcal{C}\!\mathit{omplexes}}
% \Pic is the operator that assigns to X its picard group, usage \Pic(X)
% \Picardstack_{X/B} denotes the Picard stack of X over B
% \Picardfunctor_{X/B} denotes the Picard functor of X over B
\def\Pic{\mathop{\mathrm{Pic}}\nolimits}
\def\Picardstack{\mathcal{P}\!\mathit{ic}}
\def\Picardfunctor{\mathrm{Pic}}
\def\Deformationcategory{\mathcal{D}\!\mathit{ef}}
\usepackage{cleveref}
\crefname{originaltheorem}{original theorem}{original theorems}
\Crefname{originaltheorem}{Original theorem}{Original theorems}
\makeatletter
\newlabel{etale-cohomology-theorem-topological-invariance}{{04DZ}{source}{}{sourcepointer-1}{}}
\newlabel{etale-cohomology-theorem-topological-invariance@cref}{{[originaltheorem][0][]04DZ}{[1][0][]source}{}{}{}}
\makeatother
\crefname{originaldefinition}{original definition}{original definitions}
\Crefname{originaldefinition}{Original definition}{Original definitions}
\makeatletter
\newlabel{more-morphisms-definition-thickening}{{04EX}{source}{}{sourcepointer-2}{}}
\newlabel{more-morphisms-definition-thickening@cref}{{[originaldefinition][0][]04EX}{[1][0][]source}{}{}{}}
\makeatother
\crefname{originallemma}{original lemma}{original lemmas}
\Crefname{originallemma}{Original lemma}{Original lemmas}
\makeatletter
\newlabel{morphisms-lemma-generalizations-lift-flat}{{03HV}{source}{}{sourcepointer-3}{}}
\newlabel{morphisms-lemma-generalizations-lift-flat@cref}{{[originallemma][0][]03HV}{[1][0][]source}{}{}{}}
\makeatother
\crefname{originaldefinition}{original definition}{original definitions}
\Crefname{originaldefinition}{Original definition}{Original definitions}
\makeatletter
\newlabel{topologies-definition-big-etale-site}{{021A}{source}{}{sourcepointer-4}{}}
\newlabel{topologies-definition-big-etale-site@cref}{{[originaldefinition][0][]021A}{[1][0][]source}{}{}{}}
\makeatother
\title{A sheaf of abelian groups on the big etale site vanishing below dimension n, nonzero from dimension n, invariant under thickenings}
\author{The Stacks Project Authors}
\date{stacks-project:examples.tex:local-snapshot-20261012; GNU FDL 1.2 or later}
\begin{document}
\maketitle
\noindent\textit{Editorial note.} All delivered mathematical text is the author's own, copied from the pinned Stacks Project chapter file examples.tex; no statement or proof step was rewritten, reordered or removed.\par\medskip
\noindent\textit{Editorial note.} The counted proof is the author's complete argument: The counted proof is 'See the discussion above'; the context carries the complete n=1 verification and the G\_n variant is derived by the author's explicit one-shows analogy step - recorded as mechanical analogy, content verbatim.\par\medskip
\noindent\textit{Editorial note.} The label lemma-sheaf-zero-on-low-dimension names the counted statement, and the proof environment immediately following it in the source is delivered verbatim as its proof.\par\medskip
\noindent\textit{Editorial note.} The counted statement is carried by the author's own lemma environment in the chapter source.\par\medskip
\noindent\textit{Allowed standard prerequisites (editorial).}\begin{itemize}
\item Context span 4517-4604: F(X) = product over pairs of distinct points with x' leading to x of Z/2Z, the contravariant F(f), sheafification G; the three verifications - G(X) != 0 when a nontrivial specialization exists (generalizations lift along flat morphisms, Morphisms lemma-generalizations-lift-flat), G(X) = 0 when no specializations (fields and more generally rings with all primes maximal), G(X) = G(X') for thickenings (topological invariance of the etale site, Etale Cohomology theorem-topological-invariance); the variant F\_n on |X|\textasciicircum{}\{n+1\} with sheafification G\_n
\item More on Morphisms, definition-thickening
\end{itemize}\medskip
\ifdefined\hypertarget\hypertarget{sourcepointer-1}{}\fi
\noindent\textit{Source reference.} Original theorem 04DZ: sources/stacks-chow/tags.txt, source lines 5745--5745. Prerequisite pointer invoked at 4693; it names a standard fact of a sibling chapter that the counted argument uses, and every step of the counted proof built on it is retained verbatim in the delivered body.\par\medskip
\ifdefined\hypertarget\hypertarget{sourcepointer-2}{}\fi
\noindent\textit{Source reference.} Original definition 04EX: sources/stacks-chow/tags.txt, source lines 5778--5778. Prerequisite pointer invoked at 4688; it names a standard fact of a sibling chapter that the counted argument uses, and every step of the counted proof built on it is retained verbatim in the delivered body.\par\medskip
\ifdefined\hypertarget\hypertarget{sourcepointer-3}{}\fi
\noindent\textit{Source reference.} Original lemma 03HV: sources/stacks-chow/tags.txt, source lines 4581--4581. Prerequisite pointer invoked at 4674; it names a standard fact of a sibling chapter that the counted argument uses, and every step of the counted proof built on it is retained verbatim in the delivered body.\par\medskip
\ifdefined\hypertarget\hypertarget{sourcepointer-4}{}\fi
\noindent\textit{Source reference.} Original definition 021A: sources/stacks-chow/tags.txt, source lines 2663--2663. Prerequisite pointer invoked at 4622; it names a standard fact of a sibling chapter that the counted argument uses, and every step of the counted proof built on it is retained verbatim in the delivered body.\par\medskip
\noindent
In the following we fix a big \'etale site $\Sch_\etale$
as constructed in
Topologies, Definition \ref{topologies-definition-big-etale-site}.
Moreover, a scheme will be an object of this site.
Recall that if $x, x'$ are points of a scheme $X$ we say $x$ is a
{\it specialization} of $x'$ or we write $x' \leadsto x$ if
$x \in \overline{\{x'\}}$. This is true in particular if $x = x'$.

\medskip\noindent
Consider the functor $F : \Sch_\etale \to \textit{Ab}$
defined by the following rules:
$$
F(X) = \prod\nolimits_{x \in X}
\prod\nolimits_{x' \in X, x' \leadsto x, x' \not = x}
\mathbf{Z}/2\mathbf{Z}
$$
Given a scheme $X$ we denote $|X|$ the underlying set of points.
An element $a \in F(X)$ will be viewed as a map of sets
$|X| \times |X| \to \mathbf{Z}/2\mathbf{Z}$, $(x, x') \mapsto a(x, x')$
which is zero if $x = x'$ or if $x$ is not a specialization of $x'$.
Given a morphism of schemes $f : X \to Y$ we define
$$
F(f) : F(Y) \longrightarrow F(X)
$$
by the rule that for $b \in F(Y)$ we set
$$
F(f)(b)(x, x') =
\left\{
\begin{matrix}
0 & \text{if }x\text{ is not a specialization of }x' \\
b(f(x), f(x')) & \text{else.}
\end{matrix}
\right.
$$
Note that this really does define an element of $F(X)$. We claim that if
$f : X \to Y$ and $g : Y \to Z$ are composable morphisms then
$F(f) \circ F(g) = F(g \circ f)$. Namely, let $c \in F(Z)$ and let
$x' \leadsto x$ be a specialization of points in $X$, then
$$
F(g \circ f)(x, x') = c(g(f(x)), g(f(x'))) = F(g)(F(f)(c))(x, x')
$$
because $f(x') \leadsto f(x)$. (This also works if $f(x) = f(x')$.)

\medskip\noindent
Let $G$ be the sheafification of $F$ in the \'etale topology.

\medskip\noindent
I claim that if $X$ is a scheme and $x' \leadsto x$ is a specialization
and $x' \not = x$, then $G(X) \not = 0$. Namely, let $a \in F(X)$ be
an element such that when we think of $a$ as a function
$|X| \times |X| \to \mathbf{Z}/2\mathbf{Z}$ it is nonzero at $(x, x')$.
Let $\{f_i : U_i \to X\}$ be an \'etale covering of $X$. Then we can pick an
$i$ and a point $u_i \in U_i$ with $f_i(u_i) = x$. Since generalizations
lift along flat morphisms (see
Morphisms, Lemma \ref{morphisms-lemma-generalizations-lift-flat})
we can find a specialization $u'_i \leadsto u_i$
with $f_i(u'_i) = x'$. By our construction above we see that
$F(f_i)(a) \not = 0$. Hence $a$ determines a nonzero element of $G(X)$.

\medskip\noindent
Note that if $X = \Spec(k)$ where $k$ is a field (or more generally
a ring all of whose prime ideals are maximal), then $F(X) = 0$
and for every \'etale morphism $U \to X$ we have $F(U) = 0$ because there
are no specializations between distinct points in fibres of an \'etale
morphism. Hence $G(X) = 0$.

\medskip\noindent
Suppose that $X \subset X'$ is a thickening, see
More on Morphisms, Definition \ref{more-morphisms-definition-thickening}.
Then the category of schemes \'etale over $X'$ is equivalent to the
category of schemes \'etale over $X$ by the base change functor
$U' \mapsto U = U' \times_{X'} X$, see
\'Etale Cohomology,
Theorem \ref{etale-cohomology-theorem-topological-invariance}.
Since it is always the case that $F(U) = F(U')$ in this situation
we see that also $G(X) = G(X')$.

\medskip\noindent
As a variant we can consider the presheaf $F_n$ which associates
to a scheme $X$ the collection of maps
$a : |X|^{n + 1} \to \mathbf{Z}/2\mathbf{Z}$ where $a(x_0, \ldots, x_n)$
is nonzero only if $x_n \leadsto \ldots \leadsto x_0$ is a sequence of
specializations and $x_n \not = x_{n - 1} \not = \ldots \not = x_0$.
Let $G_n$ be the sheaf associated to $F_n$.
In exactly the same way as above one shows that $G_n$ is nonzero
if $\dim(X) \geq n$ and is zero if $\dim(X) < n$.


\begin{lemma}
\label{lemma-sheaf-zero-on-low-dimension}
There exists a sheaf of abelian groups $G$ on
$\Sch_\etale$ with the following properties
\begin{enumerate}
\item $G(X) = 0$ whenever $\dim(X) < n$,
\item $G(X)$ is not zero if $\dim(X) \geq n$, and
\item if $X \subset X'$ is a thickening, then $G(X) = G(X')$.
\end{enumerate}
\end{lemma}

\begin{proof}
See the discussion above.
\end{proof}

\end{document}
